\documentclass[11pt,a4paper]{article}
\usepackage[margin=1in]{geometry}
\usepackage{hyperref}
\usepackage{enumitem}
\usepackage{graphicx}
\usepackage{xcolor}

% -----------------------------------------------------
% bvraghav-tiet-ucs503p-article.sty
% -----------------------------------------------------

% Variable: Lab Instructor
\makeatletter \newcommand{\BvrSetLabInstructor}[1]{%
  \gdef\Bvr@LabInstructor{#1}}
% default
\newcommand{\Bvr@LabInstructor}{Dr. Raghav %
  B. Venkataramaiyer}
% public accessor
\newcommand{\BvrLabInstructorValue}{\Bvr@LabInstructor}
\makeatother

% Add subtitle
\usepackage{titling}
% -- Set title bold
\pretitle{\begin{center}\bfseries\fontsize{18bp}{18bp}\selectfont}
\makeatletter
\newcommand{\BvrSubtitle}[1]{%
  \posttitle{%
    \par\end{center}
    \begin{center}\large#1\end{center}
    \vskip 0.5em}%
}
\makeatother

% Hints in the section title(s)
\newcommand{\BvrHint}[1]{%
  {\normalfont\normalsize\itshape\hfill #1}}

% -----------------------------------------------------

\title{CampusXchange: A Campus-Verified Marketplace\\ for Selling, Lending, and Exchange}

\BvrSetLabInstructor{Anushka Ma'am}

\BvrSubtitle{%
  Project Proposal for UCS503P  \\
  Submitted to: \BvrLabInstructorValue \\ \bfseries
  Thapar Institute of Engineering and Technology%
}

\author{
  Garav Goyal, CSED
  \and
  Parth Zutshi, CSED
  \and
  Dhruv Rajput, CSED
  \and
  Harjap Singh Benipal, CSED
}

\date{19 August 2026}

\begin{document}
\maketitle

\section{Higher-order goal%
  \BvrHint{\ldots secondary framing}}
This project aims to improve trust-based resource sharing within a campus
community by reducing the friction, uncertainty, and safety gaps in how
students currently buy, sell, lend, and exchange everyday items. While the
underlying need is broader than any single campus, the immediate engineering
objective is to translate \emph{safe, verified peer-to-peer commerce} into a
measurable, deployable software solution.

\section{Time-to-value %
\BvrHint{\ldots primary heuristic}}
\begin{itemize}[leftmargin=0.7cm]
    \item We will deliver meaningful increments quickly by building the Sell module first (the simplest complete transaction loop) and validating it with real users within a short iteration window.
    \item We will prioritize fast feedback over perfection by using weekly iterations (and targeting CI-backed automation early) to reduce release friction.
    \item Success is defined by what can be delivered and measured in the first iteration, then improved based on user feedback.
\end{itemize}

\section{Problem Statement}
Campus commerce today runs on ad hoc, trust-optional channels — hostel
WhatsApp groups, batch chats, and a noticeboard here and there — that were
never designed for hundreds of people who mostly don't know each other to
buy, sell, lend, or swap items. The recurring pain points are:
\begin{itemize}[leftmargin=0.7cm]
    \item \textbf{No single place to look:} a listing might live in one hostel-wing group, one batch chat, or a physical board, so it rarely reaches anyone outside the poster's own circle — buyers and sellers who would gladly trade simply never find each other.
    \item \textbf{Strangers, no proof:} anyone can be added to these groups, including people with no real link to the college. There is no reliable way to confirm the other party is even a current student before money or an item changes hands.
    \item \textbf{Lending has no memory:} every exchange is treated as a one-off sale. Nothing tracks who borrowed what, when it's due back, or whether a deposit was returned — so lending barely happens at all, even though it's the natural fit for short-term-use items like calculators, lab coats, and textbooks.
    \item \textbf{No recourse when it goes wrong:} if an item arrives damaged, a seller vanishes after payment, or a borrowed item never comes back, there is no escrow, no verified handoff record, and no one to appeal to.
\end{itemize}

\textbf{Why this matters now (contextual relevance):} A campus effectively re-runs the same shopping list every semester — the same textbooks, calculators, and lab coats change hands within a fixed, repeatedly-renewing population. That repetition means even a modest fix to trust and discoverability compounds across thousands of small trades a year, in a way it wouldn't for a marketplace serving one-off strangers.

\section{Proposed Solution %
\BvrHint{\ldots Software Proposition}}
\subsection{Overview}
Build a campus-verified mobile marketplace app, ``CampusXchange'', with the
following components:
\begin{itemize}[leftmargin=0.7cm]
    \item \textbf{Verified access:} sign-up gated by a college email address (must end in \texttt{@thapar.edu}) plus each student's Thapar ID (roll number), entered by hand and checked for a valid format — no physical ID scanning or third-party document-verification vendor involved.
    \item \textbf{Sell, Lend/Borrow, and Exchange modules:} photo-first listings with price, condition, or rent terms; category-matched exchange requests.
    \item \textbf{In-app chat and negotiation:} real-time messaging and offer cards between buyer/borrower and seller/lender.
    \item \textbf{In-app payments:} a shared, escrow-style transaction core used by all three modules.
    \item \textbf{QR-verified handoff:} a dynamic, transaction-scoped QR code confirms the physical meetup before funds are released.
    \item \textbf{Trust and safety layer:} AI-assisted content moderation with mandatory human review, plus admin-adjudicated dispute handling.
\end{itemize}

\subsection{Core workflow}
\begin{enumerate}[leftmargin=0.7cm]
    \item Student creates a listing (sell, lend, or exchange) and receives a listing ID immediately.
    \item System checks completeness and routes the listing to a moderation queue for human review before it goes live.
    \item Interested students discover the listing, negotiate over in-app chat, and agree on terms (price, or rent duration and deposit for lend/borrow).
    \item Both parties confirm the transaction; payment is held in escrow and a dynamic, time-bound QR code is issued.
    \item At the in-person handoff, the QR is scanned and condition photos are captured; escrow is released, or the transaction is flagged for admin review if a discrepancy is reported.
\end{enumerate}

\subsection{Operational constraints for a single-campus pilot setting}
\begin{itemize}[leftmargin=0.7cm]
    \item Scope the pilot to one college end-to-end: sign-up accepts only the \texttt{@thapar.edu} domain, meeting-location suggestions are pre-seeded with on-campus landmarks (hostels, academic blocks, main gate), and there is no cross-campus listing visibility at MVP.
    \item Design around semester-driven load rather than steady traffic: expect sharp spikes at the start of semester (Sell/Exchange for used textbooks and lab gear) and around move-out (graduating seniors clearing belongings), with quieter stretches mid-semester.
    \item Build for the phones and networks students actually have: a lightweight cross-platform app that stays usable on budget Android devices and patchy hostel/campus Wi-Fi, rather than assuming high-end hardware or constant connectivity.
    \item Keep every ``smart'' piece — moderation flagging, exchange matching — heuristic- and API-assisted rather than a custom-trained model, so effort goes into transaction correctness and workflow, not ML research.
    \item Meetups stay physical and on-campus by design, with no shipping or courier integration — which means the QR handoff and meeting-point selection are the entire physical-safety surface the project needs to get right.
\end{itemize}

\section{Solution Approach %
\BvrHint{\ldots Engineering focus}}
This is an engineering course project, so we focus on scalable marketplace
workflow and trust-and-safety mechanisms rather than novel algorithm
research.

\subsection{Moderation and verification assistance %
\BvrHint{\ldots non-research}}
Content moderation and identity checks will lean on simple rules and
ready-made external services, not any custom-built AI model:
\begin{itemize}[leftmargin=0.7cm]
    \item Use an off-the-shelf vision/text moderation API to assign listings a confidence-scored flag, without claiming state-of-the-art detection.
    \item Verify identity using the student's Thapar ID (roll number) entered at sign-up: check that it is a valid format, and, where a registrar-provided student list is available for the pilot, cross-check it against that list — no barcode scanning or OCR of a physical ID card is required.
    \item Show flagged listings and disputed transactions to a human moderator for decision; never fully auto-reject or auto-ban.
\end{itemize}

\subsection{Mobile and backend architecture %
\BvrHint{\ldots mobile-based solution}}
A typical mobile-plus-API architecture:
\begin{itemize}[leftmargin=0.7cm]
    \item \textbf{Frontend:} a cross-platform mobile app (listings feed, chat, transaction, and profile screens) with a separate web dashboard for admins.
    \item \textbf{Backend API:} authentication, listing management (create/edit/remove), matching/notifications, chat, and the shared transaction (escrow and QR) core.
    \item \textbf{Database:} store users, listings, chats, transactions, escrow holds, QR tokens, condition photos, and disputes.
\end{itemize}

\subsection{CI/CD and fast delivery enabler}
The project already commits, in Section 2, to CI-backed automation as part
of its time-to-value heuristic — with a full-stack mobile client, backend
API, and a payments/escrow path in scope, that automation earns its place
rather than being boilerplate. CI/CD will be used to:
\begin{itemize}[leftmargin=0.7cm]
    \item Automatically build and run tests on every change.
    \item Produce repeatable deployment artifacts to staging.
    \item Enable safe and frequent releases of incremental features.
\end{itemize}

\section{Evaluation Criterion %
\BvrHint{\ldots measurable and attributable}}
We define success using measurable metrics tied to the core workflow, chosen
so each one traces back to a pain point from Section 3.

\subsection{Primary evaluation metric}
\textbf{Time-to-listing-clearance (TLC):} time between a listing's
submission and its moderation decision (published or rejected) appearing in
the system — a check that verified, moderated access (Section 3: ``strangers,
no proof'') doesn't become a new bottleneck of its own.
\begin{itemize}[leftmargin=0.7cm]
    \item Target: median TLC $\le$ 4 hours during pilot window.
    \item Attribution: measured from database timestamps (submission vs. moderation decision by staff).
\end{itemize}

\subsection{Secondary evaluation metrics}
\begin{itemize}[leftmargin=0.7cm]
    \item \textbf{Transaction completion rate:} percentage of agreed transactions that reach a completed, QR-verified handoff — tests whether the escrow-and-QR safety net (``no recourse when it goes wrong'') actually closes the loop instead of stalling.
    \item \textbf{Lend/borrow share and on-time return rate:} share of transactions that are lend/borrow rather than sell, and the percentage of those returned by the agreed date without a dispute — tests whether native lending support (``lending has no memory'') gets used, not just built.
    \item \textbf{Dispute rate:} percentage of lend/borrow transactions flagged for admin review after condition-photo comparison.
    \item \textbf{Discovery over off-platform channels:} share of pilot users who report, via a one-question check-in, finding or listing an item in-app rather than falling back to WhatsApp or noticeboards — a proxy for whether fragmentation (``no single place to look'') is actually going down.
    \item \textbf{User clarity:} buyer/seller-reported clarity of transaction status using a simple 1--5 feedback question.
    \item \textbf{Reliability:} availability of staging/production for login, listing, and payment during peak hours (e.g., $\ge$ 99\% uptime in pilot).
\end{itemize}

\subsection{Pilot validation plan %
\BvrHint{\ldots high-level}}
\begin{itemize}[leftmargin=0.7cm]
    \item Recruit 10--20 pilot users across freshers, graduating seniors, and hostel residents, plus 2--3 volunteer moderators (or a simulated admin role if real staff are unavailable).
    \item Compare metrics before/after within the iteration window by logging baseline estimates via short interviews on current informal buying, selling, and lending habits.
\end{itemize}

\section{Scalability %
\BvrHint{\ldots with foresight}}
The proposition should scale in both theoretical and practical senses.

\begin{enumerate}[leftmargin=0.7cm]
    \item \textbf{Scalable (theoretically):} The system should support growth in number of listings and concurrent chats/transactions by:
    \begin{itemize}[leftmargin=0.7cm]
        \item decoupling components (mobile client/backend),
        \item enabling pagination and indexing for listing and category search,
        \item using stateless API design for the transaction core.
    \end{itemize}

    \item \textbf{Operationally deployable (practically):} The system should run on common cloud hosting:
    \begin{itemize}[leftmargin=0.7cm]
        \item managed database,
        \item containerized or platform-hosted deployment,
        \item CI/CD pipeline for staging/production.
    \end{itemize}
\end{enumerate}

\section{Engine availability heuristic}
A mature set of ``engines'' (tooling/services) is available to implement
this as a mobile-based project:
\begin{itemize}[leftmargin=0.7cm]
    \item Cross-platform mobile framework, plus a REST/WebSocket API layer.
    \item Managed relational database, with an in-memory cache for sessions.
    \item Authentication approach (email OTP, token-based sessions) and simple roll-number/registrar-list cross-checking — no third-party KYC or document-verification vendor required for the pilot.
    \item Object storage for listing images and condition photos.
    \item Payment gateway with escrow/split-payment support.
    \item Test framework for unit/integration tests; CI runner and staging deployment target.
\end{itemize}
We will use these mature components to focus on marketplace workflow
design, trust-and-safety correctness, and measurement rather than
inventing new infrastructure.

\section{Project scope and deliverables %
\BvrHint{\ldots iteration-friendly}}
\subsection{Initial deliverable %
\BvrHint{\ldots first iteration}}
\begin{itemize}[leftmargin=0.7cm]
    \item Student auth (college email OTP + Thapar roll-number check) and basic Sell-listing management (create, edit, remove)
    \item Buyer chat/negotiation and seller-picks-buyer flow
    \item Basic in-app payment confirmation, with escrow logic stubbed initially
    \item Basic logging and timestamps for TLC measurement
    \item CI: automatic build and test runs on every code change
    \item CD to staging with a single command or pipeline job
\end{itemize}

\subsection{Subsequent deliverables}
\begin{itemize}[leftmargin=0.7cm]
    \item Lend/Borrow and Exchange modules
    \item QR-verified handoff and mandatory condition-photo capture
    \item AI-assisted moderation queue with human review
    \item Admin dashboard for metrics, moderation, and dispute resolution
    \item Improved search/filtering and indexed category queries for scale readiness
\end{itemize}

\section{Risks and mitigations}
\begin{itemize}[leftmargin=0.7cm]
    \item \textbf{Cold-start / liquidity:} seed the marketplace with orientation-week campaigns and incentives for early listers.
    \item \textbf{Off-platform leakage:} keep the in-app QR handoff simple and fast so there is little reason to move the transaction elsewhere.
    \item \textbf{Data privacy / consent:} minimize collected identity data — roll number only, no ID card image or scan; encrypt condition-photo storage; store only what is needed.
    \item \textbf{Dispute and damage adjudication ambiguity:} instrument mandatory before/after condition photos and define clear escalation and penalty tiers before development begins.
    \item \textbf{Operational issues in pilot:} use staging early; ensure CI/CD produces repeatable deployments.
\end{itemize}

\section{Summary}
This project proposes a deployable mobile marketplace to reduce
fragmentation and safety gaps in campus buying, selling, and lending. It
meets the course criteria by emphasizing time-to-value through rapid
iterations, implementing CI/CD to support frequent safe releases, and
using measurable evaluation criteria (especially time-to-listing-clearance).
It remains focused on engineering workflow, trust-and-safety design, and
scalability readiness rather than research-driven novelty.

\end{document}
